\subsection*{Paso 3. Identifique las carencias de su infraestructura actual}

Como resultado de la evaluación técnica realizada en el Paso 2, se determina que la infraestructura diseñada se alinea completamente con los principios de una red moderna, tal como se definen en el Paso 1.

Las carencias estructurales habitualmente asociadas a redes heredadas o no planificadas (tales como dominios de broadcast extensos, topologías planas que comprometen la seguridad, o una gestión de direccionamiento IP deficiente) han sido solventadas de manera proactiva mediante el diseño implementado.

El modelo Jerárquico asegura la modularidad y la predictibilidad del rendimiento. La segmentación de VLAN y el direccionamiento IP ordenado resuelven los requisitos de aislamiento de tráfico, seguridad y gestión.

Por consiguiente, se concluye que la arquitectura evaluada cumple con los requisitos funcionales y de diseño del estándar moderno, y no presenta brechas de diseño ni deficiencias funcionales estructurales que requieran una modernización correctiva.